\documentclass[11pt]{article}
\usepackage[a4paper,margin=25mm]{geometry}
\usepackage{amsmath,amssymb,amsthm,mathtools}
\usepackage{booktabs,array,microtype,xurl}
\usepackage[hidelinks]{hyperref}

\newtheorem{theorem}{Theorem}
\newtheorem{lemma}[theorem]{Lemma}
\newtheorem{proposition}[theorem]{Proposition}
\newtheorem{corollary}[theorem]{Corollary}
\newtheorem{conjecture}[theorem]{Conjecture}
\theoremstyle{remark}
\newtheorem{remark}[theorem]{Remark}

\newcommand{\CU}{\operatorname{CatchUp}}
\newcommand{\sm}[1]{\sigma(#1)}

\title{An AI Research Attempt on the Catch-Up Optimal-Draw Conjecture\\
\large Exact response certificates, local block theorems, and obstructions to induction}
\author{AI-assisted mathematical contribution (OpenAI GPT-6 Astra Pro)}
\date{28 October 2026}

\begin{document}
\maketitle

\section*{Abstract}
We study the open Catch-Up conjecture for $S_N=\{1,\ldots,N\}$: if
$N\equiv0$ or $3\pmod4$, optimal play should give an exact draw.  This note
\emph{does not prove the conjecture}.  Its contributions are fourfold.
First, an independent exact minimax rerun confirms the draw for every
admissible $N\le28$ and reproduces the previously posted we-solve outcome signs
through $N=22$; in the checked range every opening is drawing for every
admissible $N\ge12$.  Second, we prove directly that every interval of four
consecutive integers is a drawn Catch-Up game, giving a natural local
four-block decomposition of all target prefixes.  Third, we record a
self-contained block theorem for separated additive triples, a setting in
which such local blocks really do compose and which includes Fibonacci-type
families.  Fourth, triangular numbers provide a stress test: a rigorous
four-token deficit lemma and exact draws through triangular-prefix length
$35$ coexist with explicit reachable counterexamples to the natural block
closure and singleton-reduction arguments.  The combined evidence isolates
the missing step for consecutive integers: a global invariant controlling
\emph{interleaving between adjacent four-blocks}, not merely equal partitions
or locally drawn blocks.


\paragraph{Status and provenance.}
This is an AI-generated exploration prepared for the we-solve problem
``The Catch-Up optimal-draw conjecture for consecutive integers''\footnote{%
\url{https://we-solve.github.io/problems/the-catch-up-optimal-draw-conjecture-for-consecutive-integers/}}.
It combines exact calculations and mathematical arguments from the supplied
research packages with new consolidation and verification.  Finite
computation is explicitly separated from proof.

\section{Game value and the target}
Let $R$ be the remaining finite set and let $d\ge0$ be the deficit of the
player next to select.  A selected number $x<d$ leaves the same player active;
$x\ge d$ ends the turn and transfers activity.  Equality therefore transfers
the turn.  With payoff equal to the active player's final score minus the
other score, the exact minimax value satisfies
\begin{equation}\label{eq:V}
V(\varnothing,d)=-d,\qquad
V(R,d)=\max_{x\in R}
\begin{cases}
V(R\setminus\{x\},d-x),&x<d,\\[2pt]
-V(R\setminus\{x\},x-d),&x\ge d.
\end{cases}
\end{equation}
If $d\ge\sm R$, positivity gives the useful shortcut
\begin{equation}\label{eq:terminal}
V(R,d)=\sm R-d.
\end{equation}
For $S_N=\{1,\ldots,N\}$ the conjecture is
\begin{equation}\label{eq:target}
\boxed{\quad N\equiv0,3\pmod4\quad\Longrightarrow\quad V(S_N,0)=0.\quad}
\end{equation}
The total $N(N+1)/2$ is even exactly in these two residue classes.

For an opening $x$, the responder is next with deficit $x$, so
\begin{equation}\label{eq:opening}
V(S_N,0)=\max_{x\in S_N}-V(S_N\setminus\{x\},x).
\end{equation}
Consequently a particularly convenient exact certificate for a draw is
\begin{equation}\label{eq:certificate}
V(S_N\setminus\{x\},x)\ge0\quad(1\le x\le N),
\qquad
\min_xV(S_N\setminus\{x\},x)=0.
\end{equation}
This formulation is important: an equal-sum partition only proves that some
drawn allocation exists; it is not an equilibrium strategy.

\section{Direct progress for consecutive integers}

\subsection{Independent exact verification through $N=28$}
An exact threshold implementation of \eqref{eq:V} was rerun for all
$1\le N\le28$.  It reproduces the posted we-solve win/draw/loss signs through
$N=22$ and verifies
\[
V(S_N,0)=0
\quad\text{for}\quad
N=3,4,7,8,11,12,15,16,19,20,23,24,27,28.
\]
Thus every admissible $N\le28$ draws.  This is an independent rerun rather
than a historical computational record: von Stengel's 2019 seminar abstract
reports optimal-play computations through $N=34$ with the same
``draw whenever a draw is arithmetically possible'' pattern.

The opening data are more informative than the initial values.  The openings
whose exact continuation value is a draw are:
\begin{center}
\begin{tabular}{ccl}
\toprule
$N$ & $V(S_N,0)$ & Drawing openings for Player 1 \\
\midrule
3 & $0$ & $1,3$\\
4 & $0$ & $1,2$\\
7 & $0$ & $1,2,3,4$\\
8 & $0$ & all\\
11 & $0$ & all except $10$\\
12 & $0$ & all\\
15,16,19,20 & $0$ & all\\
23,24,27,28 & $0$ & all\\
\bottomrule
\end{tabular}
\end{center}
In particular, every checked admissible prefix from $N=12$ through $28$ has
\emph{every} opening drawing.  This is finite evidence only; it suggests a
stronger eventual statement but is not used as a theorem.

\subsection{A local theorem: four consecutive integers always draw}
The target prefixes have an especially simple local decomposition into
four-term intervals.  The following fact appears to be the correct local
building block.

\begin{lemma}[Four consecutive integers]\label{lem:four-consecutive}
For every integer $a\ge1$,
\[
V(\{a,a+1,a+2,a+3\},0)=0.
\]
\end{lemma}

\begin{proof}
The case $a=1$ is a direct four-token check, so assume $a\ge2$ and write
$b=a+1$, $c=a+2$, $d=a+3$.

First the second player can prevent a win after every opening.  After an
opening $a$, reply with $d$.  The original player is then behind by $3$ with
$\{b,c\}$ remaining.  Since $3\le b$, the two-token formula
\[
V(\{p,q\},h)=q-p-h\qquad(0\le h\le p<q)
\]
gives value $c-b-3=-2$.  After an opening $b$, reply with $a$ and then $d$;
the resulting lead is exactly the remaining token $c$, so the game ties.
After an opening $c$, reply with $b$ and then $d$; the resulting lead is $c$
with only $a$ left, so the opener loses by $2$.  After an opening $d$, reply
with $b$ and then $c$; the resulting lead is exactly the remaining token $a$,
so the game ties.  Hence the first player cannot force a win.

Conversely, opening $b$ prevents a loss.  If the opponent first selects $a$,
they must continue with either $c$ or $d$; the former lets the opener finish
$2$ points ahead by taking $d$, while the latter leaves deficit $c$ and the
remaining token $c$, hence a tie.  If the opponent first selects $c$, the
opener takes $d$ and ultimately wins by $2$.  If the opponent first selects
$d$, the opener takes $c$, leaves the opponent deficit $a$, and the final
$a$ produces a tie.  Therefore opening $b$ guarantees at least a draw, and
the value is exactly zero.
\end{proof}

Thus
\begin{align*}
S_{4k}&=\bigcup_{j=0}^{k-1}\{4j+1,4j+2,4j+3,4j+4\},\\
S_{4k+3}&=\{1,2,3\}\cup
\bigcup_{j=1}^{k}\{4j,4j+1,4j+2,4j+3\},
\end{align*}
and every displayed block is individually drawn from tied scores.  The
obvious next conjecture is therefore a block-composition statement for
\emph{adjacent} four-term intervals.  It cannot be justified merely by saying
that each block draws: Catch-Up permits players to interleave selections from
different blocks.

A useful warning is already visible among shifted consecutive intervals:
\[
V(\{4,5,6,7,8\},0)=2,
\]
even though the set has even total and an equal partition.  In contrast, an
auxiliary survey found no failure among intervals of lengths $4,8,12$ with
starting point $1\le a\le12$.  The natural focused subproblem is therefore
\begin{conjecture}\label{conj:interval4k}
\[
V(\{a,a+1,\ldots,a+4k-1\},0)=0\qquad(a,k\ge1).
\]
\end{conjecture}
A proof of Conjecture~\ref{conj:interval4k} would settle the
$N\equiv0\pmod4$ half of \eqref{eq:target}.

\section{A block theorem that really does compose}
The supplied work also contains a positive model showing what a successful
block induction looks like when sufficient separation is available.

\begin{theorem}[Separated additive triples]\label{thm:additive-blocks}
Let
\[
B_j=\{a_j,b_j,c_j\},\qquad 0<a_j<b_j,\qquad c_j=a_j+b_j,
\]
and assume
\[
a_{j+1}\ge a_j+2b_j\qquad(1\le j<m).
\]
Then Catch-Up on $B_1\cup\cdots\cup B_m$ is drawn under optimal play.
Opening with $c_1$ guarantees a draw.
\end{theorem}

\begin{proof}[Proof idea, with the invariant made explicit]
Put $C_j=\sum_{i<j}c_i$.  The separation condition implies
\begin{equation}\label{eq:Cj}
C_j+a_1\le a_j,
\end{equation}
and every earlier block is smaller than every later one.  The induction uses
the following losing-gap invariant.  If the remaining set is
$A\cup B_j\cup\cdots\cup B_m$ and the active deficit $g$ satisfies
\[
\sm A<g<a_j,
\]
then the active player loses.  When that player first enters some block
$B_t$, the opponent clears the earlier $c_i$'s and uses the response
\[
\begin{array}{c|c|c}
\text{entered token}&\text{response in }B_t&\text{unused token }z\\\hline
a_t&c_t&b_t\\
b_t&c_t&a_t\\
c_t&a_t,b_t&0
\end{array}
\]
to leave a new deficit $H=z+C_t+g$.  By \eqref{eq:Cj}, the preparatory
selections are legal, the residual small-token total is strictly below $H$,
and, if a suffix remains, $H<a_{t+1}$.  The same losing-gap state has therefore
been reproduced farther to the right.

To obtain the draw rather than only the no-win half, strengthen the induction
to exact-deficit states in which the deficit equals the total of a small
residual collection $A$ and every nonempty $D\subseteq A$ satisfies
$\sm D+\min D<a_1$.  Consuming $A$ gives a clean tied suffix, while entering a
block early creates the strict losing gap above.  Applying this strengthened
statement after opening $c_1$ gives a draw.  The same response table after an
arbitrary opening lets the second player prevent a loss, completing the
induction.
\end{proof}

For any increasing positive-seed Fibonacci recurrence
$u_{n+2}=u_{n+1}+u_n$, consecutive triples have the required form and the
next triple begins at $a+2b$; hence every prefix of length divisible by $3$
is drawn.  This theorem explains why the audited Fibonacci block arguments
work.  It also explains why they do not transfer directly to consecutive or
triangular numbers: the strong separation inequality is absent.

\section{Triangular numbers as a stress test for four-block induction}
Triangular prefixes provide a useful nearby laboratory because their
arithmetic is simple, their computations strongly support an analogous
draw conjecture, and several tempting induction principles fail explicitly.
Put
\[
T_k=\frac{k(k+1)}2,\qquad \mathcal T_n=\{T_1,\ldots,T_n\}.
\]

\subsection{The arithmetic part is completely settled}
\begin{proposition}\label{prop:tri-partition}
For $n\ge1$, the triangular prefix $\mathcal T_n$ has an equal-sum partition
if and only if
\[
n\ge4\qquad\text{and}\qquad n\not\equiv1\pmod4.
\]
\end{proposition}

\begin{proof}
The total is
\[
\sum_{k=1}^nT_k=\frac{n(n+1)(n+2)}6,
\]
which is odd exactly when $n\equiv1\pmod4$.  The exceptional even-total
prefixes $n=2,3$ are not partitionable because their largest terms exceed
the sums of the remaining terms.  Bases exist at lengths $4,6,7$:
\[
10=1+3+6,\quad
21+6+1=3+10+15,\quad
21+15+6=1+3+10+28.
\]
Finally
\begin{equation}\label{eq:tri-defect}
T_m+T_{m+3}=T_{m+1}+T_{m+2}+2.
\end{equation}
Given a balanced partition at length $n$, move the token $1$ to the opposite
side and add the outer pair of the next four triangular numbers to the side
that lost $1$, with the middle pair added to the other side.  The move of $1$
changes the difference by $-2$ and \eqref{eq:tri-defect} changes it by $+2$.
Thus partitionability propagates by $n\mapsto n+4$, proving all three
eligible residue classes.
\end{proof}

The strategic analogue remains unproved.  Exact response certificates verify
an optimal draw for every eligible triangular prefix through $n=35$ (24
eligible prefixes), with largest checked token $T_{35}=630$.  The run at
$n=36$ was incomplete and is not a certificate or a counterexample.

\subsection{A rigorous four-token deficit lemma}
The same constant defect in \eqref{eq:tri-defect} does yield a genuine local
game theorem.

\begin{lemma}[Four-token exact-deficit lemma]\label{lem:four-deficit}
Suppose
\[
0<h<a<b<c<d,\qquad b-a>h,\qquad a+d=b+c+h.
\]
Then
\[
V(\{a,b,c,d\},h)=0.
\]
\end{lemma}

\begin{proof}
To show $V\le0$, give the opponent a non-losing response after every first
selection.  After $a$, respond with $b$; the original mover must then choose
$c$ or $d$, and $d$ gives a draw while $c$ loses.  After $b$ or $c$, respond
with $a$ and then $d$; legality follows from $a<b-h<c-h$, and the new lead is
the sole remaining token plus $2h$, so the original mover loses.  After $d$,
respond with $b$ and then $c$; since
$b<d-h$ and $b+c=a+d-h$, the new lead equals the remaining token $a$, giving
a draw.

For $V\ge0$, if $a+b\ge d-h$, open with $d$.  Every remaining token is below
the opponent's deficit $d-h$, but every pair reaches it.  If the opponent's
two selections leave token $z$, the original mover takes $z$ and finishes
with margin
\[
d+z-(a+b+c-z)-h=2(z-a)\ge0.
\]
If instead $a+b<d-h$, open with $a$.  Replies $b$ or $c$ are met by $d$ and
draw; after reply $d$, take $b$ and then $c$, which is legal and wins.  Hence
the mover can also prevent a loss.
\end{proof}

\begin{corollary}\label{cor:tri4}
For every $m\ge1$,
\[
V(\{T_m,T_{m+1},T_{m+2},T_{m+3}\},2)=0.
\]
\end{corollary}

\begin{proof}
For $m\ge2$, use \eqref{eq:tri-defect}, $T_m>2$, and
$T_{m+1}-T_m=m+1>2$ in Lemma~\ref{lem:four-deficit}.  The case $m=1$ is a
direct four-token check.
\end{proof}

\subsection{Why the local lemma does not compose}
The triangular data give exact obstructions to the most tempting induction
steps.

\paragraph{Two drawn quartets need not remain drawn together.}
Each four-term block in Corollary~\ref{cor:tri4} draws from deficit $2$, but
\[
V(\{T_3,\ldots,T_{10}\},2)=2,
\qquad
V(\{T_4,\ldots,T_{11}\},2)=-2.
\]
Thus even adjacent locally drawn blocks can interact nontrivially through
interleaved turns.

\paragraph{A balanced block need not be strategically neutral.}
The shifted triangular block
\[
B=\{T_5,\ldots,T_{12}\}
 =\{15,21,28,36,45,55,66,78\}
\]
has an equal partition
\[
15+36+55+66=21+28+45+78=172,
\]
yet exact minimax gives
\[
V(B,0)=-6.
\]

\paragraph{The Fibonacci singleton reduction fails.}
Let
\[
H=\{T_9,\ldots,T_{16}\}
 =\{45,55,66,78,91,105,120,136\}.
\]
Exact independent calculations give
\[
V(H,0)=0,
\qquad
V(\{10\}\cup H,10)=4.
\]
Thus even when the deficit equals an available singleton and
$2\cdot10<\min H$, entering the drawn tail can be strictly better than taking
the singleton and tying.  The exact ``crumb'' reduction used in the separated
Fibonacci proof is therefore unavailable.

\paragraph{A reachable append-four closure also fails.}
Set
\[
R=\{6,10,15,28,36,55,78,105\},\qquad
Q=\{120,136,153,171\}.
\]
Exact one-selection and independent whole-turn minimax agree that
\[
\boxed{V(R,1)=0,\qquad V(R\cup Q,1)=-2.}
\]
This state is reachable legally.  In $\mathcal T_{14}$, the alternating
single-selection history
\[
1,\ 3,\ 21,\ 45,\ 91,\ 66
\]
leaves the first player to move one point behind with remaining set $R$.
Inside $\mathcal T_{18}$ the identical history leaves $R\cup Q$ instead.
Hence reachability, a drawn state, a unit deficit, and appending precisely the
next triangular quartet still do not form a valid induction hypothesis.
Neither initial prefix is a counterexample; both initial games draw.

These failures are directly relevant to the consecutive problem: they show
that ``decompose into drawn four-blocks and induct'' is not a proof unless the
induction state retains enough information to control cross-block
interleaving.

\section{What the combined evidence suggests for $S_N$}
The calculations and counterexamples rule out several broad shortcuts:
\begin{enumerate}
\item equal-sum partitionability alone does not imply an optimal draw;
\item a union of locally drawn blocks need not itself be drawn;
\item the current deficit and a drawn/losing label for the residual tail do
      not contain enough information for a general block induction;
\item Fibonacci-style separation and exact-singleton lemmas cannot simply be
      transplanted to slowly growing sequences.
\end{enumerate}

For consecutive integers, however, there is additional rigid order absent
from the triangular counterexamples.  Lemma~\ref{lem:four-consecutive} and
the finite interval survey suggest that the most focused next target is not
an arbitrary ``union of draws'' theorem but a closure theorem for
\emph{adjacent consecutive four-blocks}.  One possible strengthened induction
would classify a small family of states
\[
(\{1,\ldots,4k\}\setminus A,d)
\]
where $A$ is confined to the current boundary block and $d$ is bounded by a
small constant, and prove that every legal excursion into the next block
returns to that family.  The triangular reachable counterexamples show why
$A$ cannot be discarded from the state description.

Equivalently, the original conjecture can be attacked through the exact
response quantities
\[
A_N(x):=V(S_N\setminus\{x\},x).
\]
The required theorem is precisely
\[
A_N(x)\ge0\quad\text{for all }x,
\qquad
\min_xA_N(x)=0,
\quad N\equiv0,3\pmod4.
\]
The computation through $N=28$ gives this in every admissible case, and in
all checked admissible cases $N\ge12$ it gives the stronger identity
$A_N(x)=0$ for every $x$.  A uniform proof of either statement remains open.

\section{Computation and verification notes}
The finite calculations use no strategy restriction.  Equation~\eqref{eq:V}
is evaluated by exact memoized backward induction; the terminal shortcut
\eqref{eq:terminal} is exact.  Larger runs use Boolean threshold queries for
$V\ge0$ and $V\ge1$, which determine win/draw/loss without replacing the
underlying minimax game.  The consecutive-integer rerun agrees with the
previously posted we-solve exact table through $N=22$ and extends the
independent local check to $N=28$.

For the triangular work, a second implementation enumerating complete legal
turns was used on small states and on the quoted reachable counterexamples.
For positive deficit $d$, a nonempty selected subset $U$ can be one complete
legal turn exactly when
\[
\sum U\ge d,
\qquad
\sum U-\max U<d;
\]
at deficit zero exactly one token must be selected.  The independent
one-selection and whole-turn recurrences agree on the audited examples.
Only compact certificates are reproduced here; the large CSV tables are not
needed for the mathematical statements above.

\section{Conclusion}
The consecutive-integer conjecture remains unresolved.  The new synthesis
narrows the conceptual gap.  Consecutive target prefixes are built entirely
from four-term intervals that are each provably drawn, and exact response
certificates continue to support the conjecture.  A separate theorem shows
that block induction can be made rigorous when additive blocks have strong
separation.  Triangular numbers, however, provide explicit reachable
counterexamples showing that local draw, equal balance, exact deficit, and
adjacent four-term extension are individually insufficient closure data.

The most concrete remaining problem is therefore:
\[
\boxed{\text{Find an interval-specific invariant that survives arbitrary
interleaving of adjacent four-blocks.}}
\]
Proving Conjecture~\ref{conj:interval4k} would already settle the
$N\equiv0\pmod4$ half of the we-solve problem; extending the same invariant
to the initial three-token seed would address the $N\equiv3\pmod4$ half.

\begin{thebibliography}{9}
\bibitem{catchup}
A.~Isaksen, M.~Ismail, S.~J.~Brams, and A.~Nealen,
\emph{Catch-Up: A Game in Which the Lead Alternates},
Game \& Puzzle Design \textbf{1}(2) (2015), 38--49.
\url{https://mpra.ub.uni-muenchen.de/108784/}

\bibitem{stengel2019}
B.~von Stengel,
\emph{Computational progress on the Catch-Up game},
LSE PhD Seminar on Combinatorics, Games and Optimisation, 22 February 2019.
\url{https://www.lse.ac.uk/Mathematics/assets/documents/Events-Archive/CGO-PhD/PhD-Seminar-on-Combinatorics-2019-update.pdf}

\bibitem{wesolveproblem}
we-solve,
\emph{The Catch-Up optimal-draw conjecture for consecutive integers},
problem record Econ-14, accessed 28 October 2026.
\url{https://we-solve.github.io/problems/the-catch-up-optimal-draw-conjecture-for-consecutive-integers/}

\bibitem{wesolveA3}
we-solve research archive,
\emph{GPT 6 Astra Pro attempt on Econ-14, version 2},
Econ-14-A3, revision r1, 9 October 2026.
\url{https://we-solve.github.io/contributions/the-catch-up-optimal-draw-conjecture-for-consecutive-integers/attempt-3/r1/}
\end{thebibliography}

\end{document}
